\section{The exact original-zeta map from the Weil trace to the
completed residue
receiver}\label{the-exact-original-zeta-map-from-the-weil-trace-to-the-completed-residue-receiver}

24 September 2026. Independent derivation, RTT0--RTT11.

\subsection{RTT0. Source objects, prerequisites, and
conventions}\label{rtt0.-source-objects-prerequisites-and-conventions}

The support remains the user's \(\tau\langle Z_1;\text{no }Z_2\rangle\).
The complete arithmetic reconstruction precedes the coefficient spaces
and operations below. Addition at \(\tau\) remains retracted. No number,
parity, metric, midpoint, or additional source operation is assigned to
that support. The corpus-and-operation rule was read before this
derivation. The new workflow rule is implemented here by constructing
the actual comparison exposed by the previous distinction between
residue duality and the Weil trace.

The complete GTR0--GTR11 proof, RD0--RD9, and the relevant GIQ2, GIQ3,
GIQ8 and GIQ9 calculations were read. SDT5--SDT7 was reread for the
strong topology and bidual step. GTR10 supplies a proved topological
isomorphism, rather than a proposed self-duality: \[
\widehat\iota:\mathcal H_{\rm res}\xrightarrow{\sim}\mathcal Q'_\beta,
\qquad \mathcal Q=\mathcal B/\mathcal I,
\tag{RTT0.1}
\] where \[
\mathcal B=\{F\text{ entire}:b_{A,M}(F)=\sup_{|\Re s|\le A}
(1+|\Im s|)^M|F(s)|<\infty\text{ for every }A,M\},
\] \[
\mathcal I=\{F\in\mathcal B:F^{(j)}(\rho)=0
\text{ for every actual nontrivial zero }\rho\text{ of }\zeta,
\ 0\le j<m_\rho\}.
\tag{RTT0.2}
\] Here \(m_\rho\) is the full multiplicity. The completion in (RTT0.1)
is the completion of \(\mathcal Q_{\rm fin}\) for the residue seminorms,
not the original topology of \(\mathcal Q\). On its dense finite-support
subspace, \[
\iota h(F)=\mathcal R(F,h)
=\sum_\rho\operatorname{Res}_{s=\rho}
\frac{F(s)h(1-s)}{\zeta(s)}\,ds.
\tag{RTT0.3}
\] The denominator is original \(\zeta\). RD1 constructs the global
isolators \(e_{\rho,j}\), whose full jet at \(\rho\) is \((s-\rho)^j\)
and whose required jets at every other zero vanish. No density of their
span in the original \(\mathcal Q\) is asserted or used here.

Use the following explicit convention for the Weil form: linear in its
first argument and conjugate-linear in its second, \[
W(F,G)=\sum_\rho m_\rho F(\rho)\overline{G(1-\overline\rho)}.
\tag{RTT0.4}
\] This is GIQ's form with the two arguments exchanged. Since that form
is Hermitian, the change of convention does not change any diagonal
value. The exact maps proving its continuity and identifying its radical
are given again below in the present \(\mathcal B\) coordinate.

\subsection{RTT1. The global value map and the full strong-dual
functional}\label{rtt1.-the-global-value-map-and-the-full-strong-dual-functional}

Let \[
H=\ell^2(\mathscr Z,m),\quad
\langle x,y\rangle_+=\sum_\rho m_\rho x_\rho\overline{y_\rho},
\quad \rho^\#=1-\overline\rho,\quad (Jx)_\rho=x_{\rho^\#}.
\tag{RTT1.1}
\] The functional equation and complex conjugation preserve
multiplicities, so \(J\) is a self-adjoint unitary involution. Every
zero here is in the entire actual nontrivial divisor, not a numerical
selection.

The map \[
E:\mathcal Q\longrightarrow H,\qquad E[F]=(F(\rho))_\rho
\tag{RTT1.2}
\] is well defined and continuous. Indeed, the unconditional zero count
\(\sum_{|\Im\rho|\le T}m_\rho=O(T\log(T+2))\), already used in GIQ2,
implies \[
\|E[F]\|_+^2\le b_{1,2}(F)^2
\sum_\rho m_\rho(1+|\Im\rho|)^{-4}<\infty.
\tag{RTT1.3}
\] The corresponding map on \(\mathcal B\) kills \(\mathcal I\), so the
quotient topology makes (RTT1.2) continuous. RD1's isolator
\(e_{\rho,0}\) maps to the coordinate vector at \(\rho\). Therefore
\(E(\mathcal Q)\) contains every finitely supported value sequence and
is dense in \(H\). This proves density of the value image only; it does
not assert density of \(\mathcal Q_{\rm fin}\) in \(\mathcal Q\).

Now \[
W(F,G)=\langle EF,JEG\rangle_+.
\tag{RTT1.4}
\] In particular its full sum converges absolutely by Cauchy--Schwarz.
For a bounded set \(B\subset\mathcal Q\), put
\(c_B=\sup_{F\in B}\|EF\|_+<\infty\). Then \[
\sup_{F\in B}|W(F,G)|\le c_B\|EG\|_+.
\tag{RTT1.5}
\] Consequently \[
\mathsf W:\mathcal Q\longrightarrow\mathcal Q'_\beta,
\qquad \mathsf W(G)(F)=W(F,G)
\tag{RTT1.6}
\] is a continuous anti-linear map. Its exact kernel is \[
\mathcal N_0=\{[G]\in\mathcal Q:G(\rho)=0\text{ for every }\rho\}.
\tag{RTT1.7}
\] One containment follows from the formula; conversely testing against
\(e_{\rho,0}\) detects \(m_\rho\overline{G(\rho^\#)}\). Thus this is
exactly the higher-jet radical already present in GIQ2, not a discarded
multiplicity block.

\subsection{RTT2. A continuous original-zeta derivative multiplier on
the full
quotient}\label{rtt2.-a-continuous-original-zeta-derivative-multiplier-on-the-full-quotient}

Set \[
(\mathsf C G)(s)=\overline{G(\overline s)},\qquad
H_\zeta(s)=s^2\zeta'(1-s).
\tag{RTT2.1}
\] The first map is anti-linear, isometric for the displayed
\(b_{A,M}\), and preserves \(\mathcal I\) because complex conjugation
preserves the full original zero divisor. The second function is entire:
the only pole of \(\zeta'(1-s)\) is the double pole at \(s=0\). Since
the residue of original \(\zeta\) at 1 is 1, \[
\zeta'(1-s)=-s^{-2}+\text{a holomorphic germ at }0,
\qquad H_\zeta(0)=-1.
\tag{RTT2.2}
\] This original pole and its exact coefficient are retained.

For completeness, \(H_\zeta\) has at most polynomial growth on every
bounded real strip. An elementary exact estimate follows from
Euler--Maclaurin, with all its terms retained: \[
\zeta(z)=\frac1{z-1}+\frac12+
\sum_{k=1}^{M}\frac{B_{2k}}{(2k)!}(z)_{2k-1}
-\frac{(z)_{2M}}{(2M)!}
\int_1^\infty\widetilde B_{2M}(x)x^{-z-2M}\,dx,
\quad \Re z>1-2M.
\tag{RTT2.3}
\] Here \((z)_r=z(z+1)\cdots(z+r-1)\), with \((z)_0=1\), and
\(\widetilde B_{2M}(x)=B_{2M}(x-\lfloor x\rfloor)\). The formula follows
by applying integration by parts successively on integer intervals to
\(x^{-z}\), using \(B_j'=jB_{j-1}\), summing the boundary terms, and
taking the upper endpoint to infinity first on \(\Re z>1\). The
displayed integral converges on the larger stated half-plane and hence
gives the analytic continuation there. At the lower endpoint the even
Bernoulli boundary values give precisely the displayed rising-factorial
terms; the remainder retains its minus sign.

Choose \(M\) so that a prescribed strip satisfies
\(\Re z+2M-1\ge\epsilon>0\). Boundedness of the periodic polynomial
gives bounds for the integral and its first derivative by a constant
times \(\epsilon^{-1}\) and \(\epsilon^{-2}\), respectively; the
derivative contributes \(-\log x\) under the integral. Differentiating
(RTT2.3) therefore proves a polynomial bound for
\(\zeta'(z)+(z-1)^{-2}\) on that strip. Substitution \(z=1-s\),
multiplication by \(s^2\), and (RTT2.2) prove the assertion for
\(H_\zeta\). Thus for each \(A\) there are finite \(d_A,C_A\) with \[
b_{A,M}(H_\zeta F)\le C_A b_{A,M+d_A}(F).
\tag{RTT2.4}
\] Multiplication by \(H_\zeta\) is consequently a continuous
endomorphism of \(\mathcal B\) preserving \(\mathcal I\).

Retain the actual isolator source \[
F_0(s)=\frac{s(s-1)}8\pi^{-s/2}\Gamma(s/2)\zeta(s),
\qquad F_0(0)=F_0(1)=\frac18.
\tag{RTT2.5}
\] Its complete factors are used here to construct an inverse to
\(L[F]=[sF]\), not to replace original \(\zeta\). Define \[
(D_0F)(s)=\frac{F(s)-F_0(s)F(0)/F_0(0)}s.
\tag{RTT2.6}
\] The numerator vanishes at zero. Outside a fixed disk, division by
\(s\) has the required polynomial bound; inside the disk Cauchy's
formula bounds the removable quotient on a smaller circle by the
numerator on a larger circle. This proves
\(D_0:\mathcal B\to\mathcal B\) continuous. It preserves \(\mathcal I\),
because no nontrivial zero is zero and the numerator retains all
required vanishing orders. Moreover \[
L[D_0F]=[F],\qquad [D_0(sF)]=[F].
\tag{RTT2.7}
\] The first identity uses \(F_0\in\mathcal I\); the second uses
\((sF)(0)=0\). Thus the induced operator is the unique continuous
inverse \(L^{-1}\). No choice of isolator changes this quotient inverse.

The resulting actual primal map is \[
\boxed{\mathsf P=L^{-2}M_{H_\zeta}\mathsf C:
\mathcal Q\longrightarrow\mathcal Q.}
\tag{RTT2.8}
\] It is continuous and anti-linear. At every nontrivial zero its full
jet is exactly the jet of \[
\zeta'(1-s)\overline{G(\overline s)}.
\tag{RTT2.9}
\] Indeed the local inverse of \(L^2\) is multiplication by \(s^{-2}\),
a holomorphic unit on each such germ. Formula (RTT2.8) proves existence
on the entire quotient even though the uncorrected global expression
(RTT2.9) has a pole at zero. The double-pole correction and the explicit
inverse are not omitted.

\subsection{RTT3. The exact local rank and every nilpotent
coordinate}\label{rtt3.-the-exact-local-rank-and-every-nilpotent-coordinate}

Fix an output zero \(\sigma\), put \(\rho=1-\sigma\),
\(\eta=\overline\sigma\), and \(m=m_\sigma=m_\rho=m_\eta\). In the
retained local parameter \(t=s-\sigma\), write \[
\zeta(\rho+v)=v^m u_\rho(v),\qquad
u_\rho(0)=\frac{\zeta^{(m)}(\rho)}{m!}\ne0.
\tag{RTT3.1}
\] Direct differentiation before taking a jet gives \[
\zeta'(\rho-t)=m(-t)^{m-1}u_\rho(-t)+(-t)^m u_\rho'(-t).
\tag{RTT3.2}
\] Multiplying the entire Taylor series of \(\mathsf C G\), and reducing
only modulo the actual defining ideal \((t^m)\), gives \[
j_\sigma(\mathsf P G)
=m(-1)^{m-1}u_\rho(0)\overline{G(\eta)}\,t^{m-1}.
\tag{RTT3.3}
\] Every omitted term in this displayed quotient calculation has degree
at least \(m\); the full original expression is (RTT3.2). Thus the
anti-linear block map from the input block at \(\eta\) to the output
block at \(\overline\eta\) has rank exactly one. Its kernel is
\(t_\eta\mathbb C[t_\eta]/(t_\eta^m)\), of complex dimension \(m-1\),
and its image is the highest-jet line \(\mathbb C t_\sigma^{m-1}\).
These statements include \(m=1\): then the rank is one, the kernel is
zero, and the coefficient is the original \(\zeta'(\rho)\).

All coefficients are nonzero on their indicated one-dimensional maps.
Hence \[
\ker\mathsf P=\mathcal N_0.
\tag{RTT3.4}
\] This follows globally because membership in \(\mathcal I\) is exactly
vanishing of every full required jet. For \(m>1\) the image is in the
local socle, the annihilator of the local maximal ideal; this is a
statement about the target line, not the deletion of that maximal ideal
from the original quotient. The complete residue matrix in RD3 remains
invertible on all \(m\) jets. The rank loss belongs to \(\mathsf P\),
the specific map from the Weil trace to that full residue receiver.

\subsection{RTT4. The finite-support identity retains the original
logarithmic
derivative}\label{rtt4.-the-finite-support-identity-retains-the-original-logarithmic-derivative}

For \(G\in\mathcal Q_{\rm fin}\), (RTT3.3) shows
\(\mathsf P G\in\mathcal Q_{\rm fin}\). For arbitrary
\(F\in\mathcal Q\), use the exact germs (RTT2.9) in the original
residue: \[
\begin{aligned}
\mathcal R(F,\mathsf P G)
&=\sum_\rho\operatorname{Res}_{s=\rho}
F(s)\overline{G(1-\overline s)}\frac{\zeta'(s)}{\zeta(s)}\,ds\\
&=\sum_\rho m_\rho F(\rho)\overline{G(1-\overline\rho)}
=W(F,G).
\end{aligned}
\tag{RTT4.1}
\] The sum is finite in this section. At a multiplicity-\(m\) zero, \[
\frac{\zeta'(\rho+t)}{\zeta(\rho+t)}=
\frac m t+\frac{u_\rho'(t)}{u_\rho(t)}.
\tag{RTT4.2}
\] The second term is holomorphic and therefore contributes zero
residue. This proves why the trace keeps precisely the constant
observations and the multiplicity \(m\), while the residue pairing alone
sees the whole jet algebra. No derivative factor was set equal to one.

\subsection{RTT5. The full completed map, without a primal-density
assumption}\label{rtt5.-the-full-completed-map-without-a-primal-density-assumption}

Define \[
\boxed{\mathsf S=\widehat\iota^{-1}\mathsf W:
\mathcal Q\longrightarrow\mathcal H_{\rm res}.}
\tag{RTT5.1}
\] It is continuous and anti-linear by (RTT0.1) and (RTT1.5). It
satisfies, on its entire domain, \[
\widehat{\mathcal R}(F,\mathsf S G)=W(F,G),\qquad
\ker\mathsf S=\mathcal N_0.
\tag{RTT5.2}
\] Here the completed residue is evaluation as defined in GTR10. Formula
(RTT4.1) proves that \(\mathsf S G\) is the canonical image of
\(\mathsf P G\) in the completion whenever \(G\) is in
\(\mathcal Q_{\rm fin}\). The assertion on all of \(\mathcal Q\) follows
from the direct construction (RTT5.1), not an extension by unproved
density.

There is also an explicit convergent residue construction for every
\(G\). Direct the finite subsets \(E\subset\mathscr Z\) by inclusion and
put \[
h_E(G)=\sum_{\rho\in E}
m_\rho(-1)^{m_\rho-1}u_\rho(0)
\overline{G(\rho^\#)}\,e_{1-\rho,m_\rho-1}.
\tag{RTT5.3}
\] By RD3's exact matrix, or by substituting
\((s-\rho)\mapsto-(s-\rho)\), \[
\iota h_E(G)(F)=\sum_{\rho\in E}
m_\rho F(\rho)\overline{G(\rho^\#)}.
\tag{RTT5.4}
\] For every bounded \(B\subset\mathcal Q\), Cauchy--Schwarz gives \[
\sup_{F\in B}|W(F,G)-\iota h_E(G)(F)|
\le c_B\left(\sum_{\rho\notin E}
m_\rho|G(\rho^\#)|^2\right)^{1/2}\longrightarrow0.
\tag{RTT5.5}
\] This is convergence in the full strong dual. By (RTT0.1), \(h_E(G)\)
converges in the specified residue completion to \(\mathsf S G\). It
treats every actual zero and all multiplicities; the finite sets are an
exact convergent net, not a bounded numerical argument.

Equations (RTT3.3) and (RTT5.3) also give a precise relationship to the
primal map. If \(\mathsf P G=\mathsf P K\), then \(G-K\in\mathcal N_0\),
so \(\mathsf S G=\mathsf S K\). There is consequently a well-defined
injective linear map \[
\mathsf j_{\rm tr}:\mathsf P(\mathcal Q)\longrightarrow\mathcal H_{\rm res},
\qquad \mathsf j_{\rm tr}(\mathsf P G)=\mathsf S G.
\tag{RTT5.6}
\] Linearity follows because both \(\mathsf P\) and \(\mathsf S\) are
anti-linear. Give \(\mathsf P(\mathcal Q)\) the quotient topology
transported from \(\mathcal Q/\mathcal N_0\) by \(\mathsf P\); then
\(\mathsf j_{\rm tr}\) is continuous by (RTT5.1). This specifies its
topology. No assertion that this quotient topology is the subspace
topology from the original \(\mathcal Q\), or that every primal class
embeds in the residue completion, is needed.

\subsection{RTT6. The exact closed receiver of the
trace}\label{rtt6.-the-exact-closed-receiver-of-the-trace}

The closure of the trace image in the full strong dual is \[
\overline{\mathsf W(\mathcal Q)}^{\,\beta}
=\mathcal N_0^\perp\subset\mathcal Q'_\beta.
\tag{RTT6.1}
\] To prove this, each \(\mathsf W(G)\) annihilates \(\mathcal N_0\), so
the closure is contained in the closed annihilator. Conversely,
\(\mathsf W(\mathcal Q)\) contains each value functional
\(\delta_{\rho,0}\) up to its nonzero multiplicity scalar, by choosing
\(G=e_{\rho^\#,0}\). Let \(D_0\) be their linear span. SDT6 proves that
every continuous linear functional on \(\mathcal Q'_\beta\) is
evaluation at an actual \(F\in\mathcal Q\). If a member
\(\lambda\in\mathcal N_0^\perp\) were outside the strong closure of
\(D_0\), locally convex separation would give such an \(F\) with every
\(\delta_{\rho,0}(F)=0\) but \(\lambda(F)\ne0\). The former says
\(F\in\mathcal N_0\), contradicting the latter. This proves (RTT6.1).

Accordingly the exact closed subreceiver inside \(\mathcal H_{\rm res}\)
is \(\widehat\iota^{-1}(\mathcal N_0^\perp)\). It retains all zero
values. If any actual zero has multiplicity at least two, then
\(\mathcal N_0\ne0\), as witnessed by \(e_{\rho,1}\), and this closed
subreceiver is proper because \(\delta_{\rho,1}\) does not annihilate
that witness. This statement does not postulate the existence of a
multiple zero. It identifies exactly what the two different receivers
record for every multiplicity allowed by the original divisor.

The full positive value completion also has an actual map into the
residue receiver: \[
\mathsf A:H\longrightarrow\mathcal H_{\rm res},\qquad
(\widehat\iota\mathsf A y)(F)=\langle EF,Jy\rangle_+,
\qquad \mathsf S=\mathsf A E.
\tag{RTT6.2}
\] This map is anti-linear and continuous from the Hilbert norm to the
residue strong topology by the bound \(c_B\|y\|_+\). It is injective
because \(E(\mathcal Q)\) is dense: if all pairings with \(Jy\) vanish,
then \(Jy=0\), hence \(y=0\). In particular the trace-to-residue map
does not collapse a nonzero value direction, even though it kills the
separate higher-jet radical in its original source. For a bounded
operator \(D:H\to H\), the map \(\mathsf A D\) is therefore zero exactly
when \(D=0\). This proves an exact detection statement for the
positive-adjoint defect calculated in GTAH, without asserting a lift of
that possibly nonholomorphic multiplier to \(\mathcal Q\), or giving
\(\mathsf A(H)\) an unproved subspace-topology inverse.

\subsection{RTT7. Dilation, generator, and the actual
transfer}\label{rtt7.-dilation-generator-and-the-actual-transfer}

On the original quotient retain \[
T_a[F]=[a^sF(s)],\qquad L[F]=[sF(s)],\qquad a>0.
\tag{RTT7.1}
\] The quotient map \(\mathsf P\) commutes with \(T_a\): \(\mathsf C\)
commutes with this multiplier because \(a\) is positive real, and the
other operators in (RTT2.8) are commuting multipliers and the inverse of
\(L\). It also commutes with \(L\). These are anti-linear commutation
statements with the indicated real parameter; no complex scalar is moved
through \(\mathsf P\) without conjugation.

The continuous dual map has the different, exact action \[
\boxed{\mathsf W T_a=a(T_{a^{-1}})'\mathsf W,
\qquad \mathsf W L=(1-L')\mathsf W.}
\tag{RTT7.2}
\] For the first identity, the multiplier in \(W(F,T_aG)\) is
\(\overline{a^{1-\overline\rho}}=a^{1-\rho}\), which is exactly the
multiplier obtained by applying \(aT_{a^{-1}}\) to \(F\). For the second
it is \(\overline{1-\overline\rho}=1-\rho\). Absolute convergence
follows as before; the extra factor \(\rho\) is controlled by one
additional defining strip seminorm. Thus no differentiation of a merely
weakly convergent sum is used.

GTR10 proves \(\widehat\iota\widehat T_a=a(T_{a^{-1}})'\widehat\iota\).
Therefore (RTT7.2) is equivalent to \[
\boxed{\mathsf S T_a=\widehat T_a\mathsf S.}
\tag{RTT7.3}
\] The generator on the residue completion is explicitly
\(\widehat L=\widehat\iota^{-1}(1-L')\widehat\iota\), a continuous
operator extending the finite-support residue generator. Hence
\(\mathsf S L=\widehat L\mathsf S\). Transposition of a continuous
operator is strong-continuous because it maps each bounded set to a
bounded set; this verifies continuity of the displayed operator.

For recovered integer covers, set \(\mathsf U_n=nT_{n^{-1}}\) on actual
degree-one cohomology. GTR2--GTR4 construct it from the weighted
ramified sheaf trace and inverse coefficient action. Retaining
\(\kappa:Q_A=A/J\to\mathcal Q\),
\(\kappa[a]=[\tfrac12\int_0^\infty a(u)u^sdu/u]\), the actual receiving
map is \[
\mathsf S_A=\mathsf S\kappa:Q_A\longrightarrow\mathcal H_{\rm res},
\qquad \mathfrak I=\kappa'\widehat\iota:
\mathcal H_{\rm res}\xrightarrow{\sim}Q'_{A,\beta}.
\tag{RTT7.4}
\] It obeys \[
\mathfrak I\mathsf S_A T_n
=\mathsf U_n'\mathfrak I\mathsf S_A.
\tag{RTT7.5}
\] All domains and the Mellin factor \(1/2\) are retained. The two pole
restrictions are still \(h\mapsto(h,-h)\), so composing (RTT7.4) gives
\(G\mapsto(\mathsf S_A G,-\mathsf S_A G)\). No supported pole is
identified with the other one.

\subsection{RTT8. The geometric transfer and the Hilbert adjoint have a
proved
comparison}\label{rtt8.-the-geometric-transfer-and-the-hilbert-adjoint-have-a-proved-comparison}

The Weil form has the exact adjoint identity \[
W(T_aF,G)=W(F,aT_{a^{-1}}G).
\tag{RTT8.1}
\] Indeed the multiplier on the right is
\(\overline{a^{1-\rho^\#}}=a^\rho\), the multiplier on the left. This is
an adjoint with respect to the specified possibly indefinite and
possibly degenerate Weil form. On \(\mathcal Q/\mathcal N_0\) it is
nondegenerate; on the full \(\mathcal Q\), an adjoint identity alone
cannot determine changes whose image is in \(\mathcal N_0\). The chosen
full operator in (RTT8.1) is nevertheless already fixed by the geometric
construction and has all its higher jets.

On the positive auxiliary space \(H\), write \[
(D_ax)_\rho=a^\rho x_\rho,\qquad
(U_ax)_\rho=a^{1-\rho}x_\rho.
\tag{RTT8.2}
\] These are bounded operators, since \(0<\Re\rho<1\). Direct evaluation
in \(\langle\ ,\ \rangle_+\) gives \[
(D_a^*x)_\rho=a^{\overline\rho}x_\rho,
\qquad \boxed{U_a=J D_a^*J.}
\tag{RTT8.3}
\] The second identity follows from \(1-\rho^\#=\overline\rho\). Thus
the constructed geometric transfer receives the adjoint of pullback for
the actual Weil form, and its relation to the positive Hilbert adjoint
is the explicit conjugation (RTT8.3). The \(\zeta'\) multiplier in the
trace-to-residue map does not remove that involution; (RTT5.2) proves
that it recovers exactly the old Weil values.

This answers what additional comparison the residue bridge makes
available: any proposed change of the positive adjoint can now be
transported through \(\mathsf S\) and tested by (RTT7.2)--(RTT8.3) on
the same original quotient. One cannot replace \(U_a\) by \(D_a^*\)
merely because a residue pairing is nondegenerate. Their exact relation,
including \(J\), has been calculated here; further positivity work must
use it rather than assume that it is the identity relation.

\subsection{RTT9. Original functional-equation factors and the full
arithmetic
form}\label{rtt9.-original-functional-equation-factors-and-the-full-arithmetic-form}

No completed zeta replaces original \(\zeta\) in this construction. The
complete original functional equation is \[
\zeta(s)=\chi(s)\zeta(1-s),\qquad
\chi(s)=2^s\pi^{s-1}\sin(\pi s/2)\Gamma(1-s)
=\pi^{s-1/2}\frac{\Gamma((1-s)/2)}{\Gamma(s/2)}.
\tag{RTT9.1}
\] Its exact differentiated identity is \[
\zeta'(1-s)=\frac{\chi'(s)}{\chi(s)^2}\zeta(s)
-\frac1{\chi(s)}\zeta'(s).
\tag{RTT9.2}
\] At a nontrivial zero, \(\chi\) is a holomorphic unit. The first term
vanishes to the full original multiplicity; hence the full jet of
(RTT2.9) can also be computed from the second term, retaining its minus
sign and unit. Equation (RTT9.2), rather than a removal of those
factors, proves that local quotient comparison. Globally the definition
remains (RTT2.8), and it does not require a meromorphic multiplier
involving \(\chi^{-1}\) to preserve \(\mathcal B\).

The use of \(F_0\) in (RTT2.5) retains the endpoint values and all its
original trivial-zero values \[
F_0(-2r)=\frac{r(2r+1)(-1)^r\pi^r}{2\,r!}\zeta'(-2r)
=\frac{(1+2r)(2r)}8\pi^{-(1+2r)/2}
\Gamma((1+2r)/2)\zeta(1+2r),\quad r\ge1.
\tag{RTT9.3}
\] They are not nontrivial-zero coordinates in \(\mathcal Q\); their
role in the full source reconstruction and the four endpoint lines
remains RD8 and SSI.

More directly, (RTT5.2) is an equality with the existing complete
original arithmetic form, not a newly selected sign form. For
\(A_{F,G}(s)=F(s)\overline{G(1-\overline s)}\), choose the logarithmic
test \(h\) with \(A_{F,G}(s)=\int_{\mathbb R}h(v)e^{-(s-1/2)v}\,dv\),
whose existence in the full test space follows by the Mellin comparison
and convolution proof of GIQ9. Retain
\(\widehat h(t)=\int h(v)e^{-itv}dv\). Then \[
\widehat{\mathcal R}(F,\mathsf S G)
=A_{F,G}(0)+A_{F,G}(1)+A_\infty(h)-P_{\rm hist}(h),
\tag{RTT9.4}
\] where every original term is \[
P_{\rm hist}(h)=\sum_{n\ge2}
\frac{\log L_n-\log L_{n-1}}{\sqrt n}
\bigl(h(\log n)+h(-\log n)\bigr),
\qquad L_n=\operatorname{lcm}(1,\ldots,n),
\] \[
A_\infty(h)=\frac1{2\pi}\int_{\mathbb R}\widehat h(t)
\left(\Re\frac{\Gamma'(1/4+it/2)}{\Gamma(1/4+it/2)}-\log\pi\right)dt.
\tag{RTT9.5}
\] The coefficient \(\log L_n-\log L_{n-1}=\Lambda(n)\) retains all
prime-power repetitions from the complete history. For each finite
trivial-zero cutoff \(B\), the raw divisor identity is also retained: \[
V_{\zeta,B}=W(F,G)+\sum_{r=1}^B A_{F,G}(-2r)-A_{F,G}(1),
\] \[
G_B=A_\infty(h)+A_{F,G}(0)+\sum_{r=1}^B A_{F,G}(-2r),
\qquad V_{\zeta,B}=G_B-P_{\rm hist}(h).
\tag{RTT9.6}
\] No divergent infinite sum over trivial zeros is introduced. These
equalities follow from the fully proved GIQ9 identity and (RTT5.2). The
same support coefficients therefore carry exactly the same arithmetic
compensation before and after applying the new map.

\subsection{RTT10. Labels and the coefficient
extension}\label{rtt10.-labels-and-the-coefficient-extension}

Let \(K\) be the complex vector space of the retained formal coefficient
or support records, or their already constructed algebraic tensor
product. Because \(\mathsf P\) and \(\mathsf S\) are anti-linear, an
ordinary tensor with \(1_K\) must first have its domain typed correctly.
Write \(\overline{\mathcal Q}\) for the conjugate vector space, with \[
\lambda\cdot\overline q=\overline{\overline\lambda q}.
\] The maps \[
\widetilde{\mathsf P}:\overline{\mathcal Q}\to\mathcal Q,
\quad \overline q\mapsto\mathsf Pq,
\qquad
\widetilde{\mathsf S}:\overline{\mathcal Q}\to\mathcal H_{\rm res},
\quad \overline q\mapsto\mathsf Sq
\tag{RTT10.1}
\] are complex-linear: for example
\(\mathsf S(\overline\lambda q)=\lambda\mathsf Sq\). Thus their tensor
extensions with \(1_K\) have domain
\(\overline{\mathcal Q}\otimes_{\mathbb C}K\). Balancing holds
explicitly: \[
(\widetilde{\mathsf S}\otimes1_K)((\lambda\overline q)\otimes k)
=\lambda\mathsf Sq\otimes k
=\mathsf Sq\otimes\lambda k
=(\widetilde{\mathsf S}\otimes1_K)(\overline q\otimes\lambda k).
\tag{RTT10.2}
\] The identical argument applies to \(\widetilde{\mathsf P}\). A finite
expansion in a basis of \(K\) proves that both kernels are exactly
\(\overline{\mathcal N_0}\otimes_{\mathbb C}K\). Every record is
unchanged.

There is also an explicitly anti-linear presentation on
\(\mathcal Q\otimes_{\mathbb C}K\) if the specified formal/support basis
\((k_b)_b\) is retained. Define
\(C_K(\sum_b c_bk_b)=\sum_b\overline{c_b}k_b\); this fixes each label
and conjugates only its complex coefficient. Then \[
q\otimes k\longmapsto\mathsf Sq\otimes C_Kk
\tag{RTT10.3}
\] is well defined and anti-linear. Indeed the two representatives
\((\lambda q)\otimes k\) and \(q\otimes\lambda k\) both map to
\(\overline\lambda\mathsf Sq\otimes C_Kk\). Its kernel is
\(\mathcal N_0\otimes_{\mathbb C}K\) by the same basis argument. The
formula with \(\mathsf P\) has the same properties. This second
presentation specifies the needed coefficient conjugation instead of
implicitly identifying \(\overline{\mathcal Q}\) with \(\mathcal Q\).

Neither presentation identifies \([\tau]\) with \([1]\) or either
support branch with the other. No multiplication or addition of the
source \(\tau\) is inferred from operations on its already constructed
coefficient receiver.

\subsection{RTT11. Exact outcome and receiving
calculation}\label{rtt11.-exact-outcome-and-receiving-calculation}

There are now two proved maps with different codomains: \[
\mathsf P=L^{-2}M_{s^2\zeta'(1-s)}\mathsf C:\mathcal Q\to\mathcal Q,
\qquad
\mathsf S=\widehat\iota^{-1}\mathsf W:\mathcal Q\to\mathcal H_{\rm res}.
\tag{RTT11.1}
\] Both are continuous and anti-linear, and both have exactly the full
value-zero kernel \(\mathcal N_0\). Their finite-support comparison and
full strong limit are (RTT4.1) and (RTT5.3)--(RTT5.6). Locally the exact
original \(\zeta'\) coefficient maps each multiplicity block onto its
conjugate block's highest-jet line. Globally the completed residue
pairing of \(\mathsf S G\) is exactly the original full Weil pairing of
\(G\), including every arithmetic, Gamma, endpoint and finite
trivial-divisor contribution.

The receiving action is the actually constructed geometric transfer
transpose, (RTT7.5); the relation to the positive Hilbert adjoint is
(RTT8.3). Those maps make a further positivity calculation concrete
without changing the original Weil values or assuming purity. The
comparison itself establishes neither the absence of off-line zeros nor
a new positivity theorem.

\subsubsection{Proof-source paths actually
used}\label{proof-source-paths-actually-used}

\begin{itemize}
\tightlist
\item
  \nolinkurl{ORIGINAL_ZETA_RESIDUE_DUAL_AND_SUPPORT_MAPS.md}, RD0--RD9:
  entire quotient, global isolators, original residue matrices, full
  functional-equation units, sheaf comparison and original exceptional
  terms.
\item
  \nolinkurl{CC_RAMIFIED_TRACE_AND_RESIDUE_ACTION.md}, GTR0--GTR11:
  actual weighted trace, inverse coefficient action, strong completion,
  signed support maps and derived transfer.
\item
  \nolinkurl{GLOBAL_INFINITESIMAL_QUOTIENT.md}, GIQ2--GIQ3 and
  GIQ8--GIQ9: value form, nilpotent radical, actual dilation and
  complete original arithmetic identity.
\item
  \nolinkurl{../tau_weight_cohomology_20260924/CC_STRONG_DUAL_TOPOLOGY_AND_JETS.md},
  SDT5--SDT7: exact strong topology, strong bidual and separation
  argument. The completed residue is used through GTR10's already proved
  topological isomorphism.
\end{itemize}

No numerical test, simplicity assumption, finite-support primal-density
claim, or new publication receipt occurs in this derivation.
